\chapter[Introduction]{Introduction}{Introduction}\label{Intro}

\noindent
\rule{0.49\textwidth}{0.75pt} $_{\bigcirc}$ \rule{0.49\textwidth}{0.75pt}\\
Identity verification has been established as one of the basic prerequisites in the contemporary digital ecosystems where access to public services financial systems and private platforms are becoming increasingly dependent on government issued identity credentials. A single identity in India like Aadhaar is now the access point to Direct Benefit Transfers Aadhaar facilitated payment systems, telecom onboarding, digital lockers, passport services and various welfare schemes. This integration has helped in the increased reach of the service and efficiency of operation but on the other hand it has greatly exposed sensitive personal information. Most of the existing systems verification continues to rely on central authority and disclosure of full documents that increases the attack surface on which people can use without being detected or long term loss of privacy. A notable one is the Mukhyamantri Majhi Ladki Bahin Yojana in Maharashtra in which audits have shown that thousands of ineligible beneficiaries have been able to pass the eligibility test. 
When analyzing the failures of the public service delivery in the real world, the necessity to have a secure and privacy preserving framework of identity becomes obvious. Although the authentication on a large scale requires Aadhaar authentication, cases of massive fraud and omission still persist
The value of this work is in the fact that it introduces a pragmatic research approach that does not substitute the existing government infrastructure but acts as a supplement. TrustChain has been made to be interoperable with legitimate government services and private bodies to provide initial authenticity with privacy being implemented at the verification level. In accordance with practical constraints imposed by laws and failures in implementation practices that have been observed, this study shows that decentralized identity and ZKP based verification can be used to mitigate fraud, eliminate exclusion and protect the privacy of citizens. The proposed system will help to sustain a safer and more resilient and user controlled identity ecosystem that can be deployed on a large scale of public services
.\\
\noindent

\section{Motivation / Objective}
TrustChain is inspired by the fact that a user has been exposed to personal identity data many times through digital services, and in the course of a simple verification, users are required to provide complete documents. Cases such as welfare fraud and unsafe KYC practices indicate that current systems display more information than needed. This project aims at developing a decentralized identity verification system that can be used to establish genuineness without revealing personal information. TrustChain provides user controlled access, privacy and secure verification to real world applications by utilizing blockchain, IPFS, SSI and ZKP. 

\section{Proposed Solution}
The solution proposed is TrustChain, which is an identity verification system based on blockchain, requiring no sharing of raw identity documents. Authenticated files are uploaded by the users to IPFS once, and are represented in the form of cryptographic hash values. The system enables users to demonstrate their authenticity with Self Sovereign Identity and ZK-STARK based proofs, without exposing personal information about themselves. Smart contracts are used to offer role based access control and to process verification requests. The users can add and remove access at any time, guaranteeing complete ownership, privacy and secure authentication by institutions.

\section{Scope of the Project}
The TrustChain can be applied in various areas where identity validation is needed such as in government services, telecom, banking, education and enterprise systems. The project aims at delivering a privacy saving verification layer that is capable of integrating into existing infrastructures such as DigiLocker and Aadhaar based services without substituting them. It also works with secure onboarding and decentralized storage with the help of and verification with ZK-STARK proofs.s. The system allows role based access control and user driven consent management and revocation. Future scope also involves connecting to government databases to check authenticity, creating a simplified in built wallet to be used by non technical users and scaling the platform to be used by large population.

\section{Report Overview}
This report involves the design and development of TrustChain, a decentralized identity verification platform that puts emphasis on privacy preserving access control. It discusses the issue of too much identity exposures within the current systems and the necessity to have a secure alternative. The report outlines the architecture of the system with blockchain integration, IPFS based storage, Self Sovereign Identity and ZK-STARK based verification. It describes the entire process of registering users, verifying their identities and managing access. Other sections of the document are technology stack, implementation strategy, case study and comparisons with other solutions in place such as DigiLocker. On the whole, it offers a clear insight of how TrustChain can guarantee secure, user controlled and non intrusive identity verification. It also brings to the fore the security considerations, limitations and challenges encountered during the implementation as well as potential improvements. The report ends with the possible effects of TrustChain to the real world conditions where privacy and secured identity checks are of essence.

\rule{0.49\textwidth}{0.75pt} $_{\bigcirc}$ \rule{0.49\textwidth}{0.75pt}\\
